\section{Conjugated Systems} % lessons 1-3

  \subsection{Organic Acids \& Bases Review}

    \begin{core}[label=core:general_acidity]{General Organic Compound Acidity}

      Memorizing the following pK$_\text{a}$ values as baseline references is essential---they serve as anchors for estimating the acidities of related compounds by comparing functional groups, electron-withdrawing effects, and resonance stabilization of the conjugate base.

      \begin{table}[H]
        \centering
        \begin{tabular}{>{\raggedright\arraybackslash}p{1.9cm}cc >{\raggedright\arraybackslash}p{1.9cm}cc}
          \hline
          \hline
          \textbf{Compound} & \textbf{pK$_\text{a}$} & \textbf{Structure} &
          \textbf{Compound} & \textbf{pK$_\text{a}$} & \textbf{Structure} \\
          \hline
          \rule{0pt}{5ex}{\small Alkane} &
          {\small 50} &
          \scalebox{0.75}{
            \scheme{\chemfig[atom sep=2em]{
                [:0]\textcolor{Red}{H}
                -[:30,,,,Red]
                -[:-30]
                -[:30]
                -[:-30]
            }}
          } &
          {\small Ammonium} &
          {\small 10} &
          \scalebox{0.75}{
            \scheme{\chemfig[atom sep=2em]{
                [:0]\textcolor{Red}{H}
                -[:0,,,,Red]\pcharge{135}{N}
                (-[:45]R)
                (-[:-45]R)
                -[:0]R
            }}
          }
          \\[4ex]

          {\small Alkene} &
          {\small 45} &
          \scalebox{0.75}{
            \scheme{\chemfig[atom sep=2em]{
                [:0]\textcolor{Red}{H}
                -[:30,,,,Red]
                =[:-30]
                -[:30]
                -[:-30]
            }}
          } &
          {\small $\alpha$ to Two Carbonyls} &
          {\small 10} &
          \scalebox{0.75}{
            \scheme{\chemfig[atom sep=2em]{
                [:0]R
                -[:30]
                (=[:90]O)
                -[:-30]
                (-[:-90,,,,Red]\textcolor{Red}{H})
                -[:30]
                (=[:90]O)
                -[:-30]R
            }}
          }
          \\[4ex]

          {\small Allylic} &
          {\small 40} &
          \scalebox{0.75}{
            \scheme{\chemfig[atom sep=2em]{
                [:0]\textcolor{Red}{H}
                -[:30,,,,Red]
                -[:-30]
                =[:30]
                -[:-30]
                -[:30]
            }}
          } &
          {\small Carboxylic Acid} &
          {\small 5} &
          \scalebox{0.75}{
            \scheme{\chemfig[atom sep=2em]{
                [:0]\textcolor{Red}{H}
                -[:-30,,,,Red]O
                -[:30]
                (=[:90]O)
                -[:-30]R
            }}
          }
          \\[4ex]

          {\small Amine} &
          {\small 35} &
          \scalebox{0.75}{
            \scheme{\chemfig[atom sep=2em]{
                [:0]\textcolor{Red}{H}
                -[:0,,,,Red]N
                (-[:45]R)
                (-[:-45]R)
            }}
          } &
          {\small Oxonium} &
          {\small 0} &
          \scalebox{0.75}{
            \scheme{\chemfig[atom sep=2em]{
                [:0]\textcolor{Red}{H}
                -[:0,,,,Red]\pcharge{135}{O}
                (-[:45]R)
                (-[:-45]R)
            }}
          }
          \\[4ex]

          {\small Alkyne} &
          {\small 25} &
          \scalebox{0.75}{
            \scheme{\chemfig[atom sep=2em]{
                [:0]\textcolor{Red}{H}
                -[:0,,,,Red]C
                ~[:0]C
                -[:0]R
            }}
          } &
          {\small Trifluoro Acetic Acid} &
          {\small 0} &
          \scalebox{0.75}{
            \scheme{\chemfig[atom sep=2em]{
                [:0]\textcolor{Red}{H}
                -[:-30,,,,Red]O
                -[:30]
                (=[:90]O)
                -[:-30]\ce{CF3}
            }}
          }
          \\[4ex]

          {\small $\alpha$ to an Ester} &
          {\small 25} &
          \scalebox{0.75}{
            \scheme{\chemfig[atom sep=2em]{
                [:0]R
                -[:-30]
                (-[:-90,,,,Red]\textcolor{Red}{H})
                -[:30]
                (=[:90]O)
                -[:-30]OR
            }}
          } &
          {\small Hydrohalic Acid} &
          {\small <0} &
          \scalebox{0.75}{
            \scheme{\chemfig[atom sep=2em]{
                [:0]\textcolor{Red}{H}
                -[:0,,,,Red]X
            }}
          }
          \\[4ex]

          {\small $\alpha$ to a Ketone or Aldehyde} &
          {\small 20} &
          \scalebox{0.75}{
            \scheme{\chemfig[atom sep=2em]{
                [:0]R
                -[:-30]
                (-[:-90,,,,Red]\textcolor{Red}{H})
                -[:30]
                (=[:90]O)
                -[:-30]R
            }}
          } &
          {\small Sulfuric Acid} &
          {\small <0} &
          \scalebox{0.75}{
            \scheme{\chemfig[atom sep=2em]{
                [:0]\textcolor{Red}{H}
                -[:0,,,,Red]\ce{OSO3H}
            }}
          }
          \\[4ex]

          {\small Alcohol (ROH)} &
          {\small 15} &
          \scalebox{0.75}{
            \scheme{\chemfig[atom sep=2em]{
                [:0]\textcolor{Red}{H}
                -[:0,,,,Red]OR
            }}
          } &
          {\small Toluenesulfonic Acid} &
          {\small <0} &
          \scalebox{0.75}{
            \scheme{\chemfig[atom sep=2em]{
                [:0]\textcolor{Red}{H}
                -[:0,,,,Red]O
                -[:0]S
                (=[:90]O)
                (=[:-90]O)
                -[:0]*6(-=-
                (-[:0]\ce{CH3})
                =-=-)
            }}
          }
          \\[4ex]

          {\small Phenol (PhOH)} &
          {\small 10} &
          \scalebox{0.75}{
            \scheme{\chemfig[atom sep=2em]{
                [:0]\textcolor{Red}{H}
                -[:0,,,,Red]O
                -[:0]*6(-=-=-=-)
            }}
          } &
          & &
          \\[4ex]

          \hline
        \end{tabular}
        \caption{Common acids and their pK$_\text{a}$ values, ordered from least to most acidic}
        \label{tab:pka_table}
      \end{table}

    \end{core}

    \begin{core}{Acidity of Alkanes, Alkenes, and Alkynes}

      We now examine the pK$_\text{a}$ values trend of alkanes, alkenes, and alkynes. The linear pattern observed can be rationalized by considering the following factors.

      \reactionfig{
        \scheme{
          \chemname{
            \chemfig{
              [:0]\textcolor{Red}{H}
              -[:30,,,,Red]\textcolor{Blue}{\charge{90:6pt=\footnotesize$\text{sp}^3$}{}}
              -[:-30]
              -[:30]
            }
          }{Alkane \\ pK$_\text{a} \approx 50$}
          \qquad \qquad
          \chemname{
            \chemfig{
              [:0]\textcolor{Red}{H}
              -[:30,,,,Red]\textcolor{Blue}{\charge{90:6pt=\footnotesize$\text{sp}^2$}{}}
              =[:-30]
              -[:30]
            }
          }{Alkene \\ pK$_\text{a} \approx 45$}
          \qquad \qquad
          \chemname{
            \chemfig{
              [:0]\textcolor{Red}{H}
              -[:0,,,,Red]\textcolor{Blue}{\charge{90:6pt=\footnotesize$\text{sp}$}{\textcolor{Black}{C}}}
              ~[:0]CR
            }
          }{Alkyne \\ pK$_\text{a} \approx 25$}
        }
      }{Alkane, alkene, and alkyne pK$_\text{a}$ values}{fig:ph1a}{0.9}

      \begin{enumerate}

        \item \textbf{Electronegativity of atom attached to the acidic H}: This factor cannot account for the trend, since the acidic hydrogen is attached to a carbon atom in all three cases, precluding any difference in electronegativity.

        \item \textbf{Inductive Effect}: This factor is likewise insufficient, as only C--H bonds are present in the vicinity of the acidic proton.

        \item \textbf{Resonance}: No resonance stabilization of the conjugate base is possible, so this factor does not explain the trend.

        \item \textbf{Charge}: All species are neutral, so this factor offers no explanation.

        \item \textbf{Hybridization}: Here a genuine difference emerges, as the acidic carbon is sp$^3$-, sp$^2$-, or sp-hybridized in the alkane, alkene, and alkyne, respectively. The lone pair generated upon deprotonation is stabilized most effectively in the sp-hybridized orbital, since its greater s character holds the electrons closer to the nucleus.

      \end{enumerate}

      When the acidity cannot be rationalized from the structure of the acid itself, it is useful to instead examine the \textbf{conjugate bases}.

      \reactionfig{
        \scheme{
          \chemname{
            \chemfig{
              [:0]\ncharge{180}{}\textcolor{Blue}{\charge{60:6pt=\footnotesize$\text{sp}^3$}{}}
              -[:-30]
              -[:30]
            }
          }{Alkyl anion \\ (\textit{carbanion})}
          \qquad \qquad
          \chemname{
            \chemfig{
              [:0]\ncharge{180}{}\textcolor{Blue}{\charge{60:6pt=\footnotesize$\text{sp}^2$}{}}
              =[:-30]
              -[:30]
            }
          }{Alkenyl anion \\ (\textit{vinyl anion})}
          \qquad \qquad
          \chemname{
            \chemfig{
              [:0]\ncharge{180}{C}\textcolor{Blue}{\charge{90:8pt=\footnotesize$\text{sp}$}{}}
              ~[:0]C
              -[:0]R
            }
          }{Alkynyl anion \\ (\textit{acetylide anion})}
        }
      }{Conjugate bases of an alkane (\textit{left}), alkene (\textit{middle}), and alkyne (\textit{right})}{fig:ph2a}{0.85}

      Among the anions shown in \autoref{fig:ph2a}, the \textbf{alkynyl anion is the weakest base} because its lone pair is the most stabilized and thus least reactive. This arises because its \ce{e^-} density is drawn more toward the carbon nucleus---away from the bonding region---making it less attractive toward \ce{H+} relative to the alkenyl and alkyl anions. The following energy diagram below illustrates this comparison.

      \pgfig{
        \draw[->] (0,0) -- (0,3.2);
        \node[] at (-0.4,1.6) {$E$};

        % --- sp3 panel ---
        \begin{scope}[xshift=0.5cm]
          \draw[black] (0.2,2.5) -- (1.2,2.5);
          \draw[black,->] (0.5,3.0) -- (0.5,2.6);
          \draw[black,->] (0.9,2.6) -- (0.9,3.0);
          \node at (0.7,0.3) {$C\,(sp^3)$};
        \end{scope}

        % --- sp2 panel ---
        \begin{scope}[xshift=2.5cm]
          \draw[black] (0.2,2.0) -- (1.2,2.0);
          \draw[black,->] (0.5,2.5) -- (0.5,2.1);
          \draw[black,->] (0.9,2.1) -- (0.9,2.5);
          \node at (0.7,0.3) {$C\,(sp^2)$};
        \end{scope}

        % --- sp panel ---
        \begin{scope}[xshift=4.5cm]
          \draw[black] (0.2,1.5) -- (1.2,1.5);
          \draw[black,->] (0.5,2.0) -- (0.5,1.6);
          \draw[black,->] (0.9,1.6) -- (0.9,2.0);
          \node at (0.7,0.3) {$C\,(sp)$};
        \end{scope}
      }{Diagram of $\text{sp}^3$, $\text{sp}^2$, and sp orbital energy levels; \textit{lone pair in the sp orbital is most stabilized because the relatively greater s character increases the probability \ce{e^-}s residing closer to the positive charge of the nucleus}}{fig:ph3a}{0.8}

    \end{core}

    \begin{core}{K$_\text{eq}$ and pK$_\text{a}$s}

      Any acid–base reaction has an equilibrium constant given by $K_\text{eq} = 10^{(\text{p}K_\text{a,conj.\,acid} - \text{p}K_\text{a,acid})}$. Using the pK$_\text{a}$ values in \autoref{tab:pka_table}, we can therefore estimate the $K_\text{eq}$ of a given acid–base reaction to a reasonable degree of accuracy, as illustrated in the following examples.

      \begin{enumerate}

        \item \textbf{Protonation of ammonia by hydronium}

        \reactionfig{
          \scheme{
            \chemname{
              \chemfig{
                [:0]@{oxygen}\pcharge{60}{O}
                (-[:135]H)
                (-[:-135]H)
                -[@{bond}:0,,,,Orange]@{hydrogen}\textcolor{Orange}{H}
              }
            }{pK$_\text{a} \approx 0$}
            \+
            \chemfig{
              [:0]@{nitrogen}\ce{NH3}
            }
            \qquad \qquad
            \arrow{<->>[][$\text{K} = 10^{(10-0)}$]}
            \qquad
            \chemname{
              \raisebox{-0.5cm}{
                \chemfig{
                  [:0]\textcolor{Orange}{H}
                  -[:0,,,,Orange]\pcharge{90}{\ce{NH3}}
                }
              }
            }{pK$_\text{a} \approx 10$}
            \raisebox{-0.5cm}{
              \+
              \chemfig{
                [:0]\ce{H2O}
              }
            }
          }
          \chemmovearrow{
            \reactionfigarrow{bond}{-80}{oxygen}{-100}{Red}
            \reactionfigarrow{nitrogen}{135}{hydrogen}{45}{Red}
          }
        }{Protonation of ammonia by hydronium}{fig:ph4a}{0.9}

        \item \textbf{Enolate formation} 

        \reactionfig{
          \scheme{
            \chemfig{
              @{ho}\ce{HO-}
            }
            \+
            \chemname{
              \chemfig{
                [:0]R
                -[:30]@{carbon}
                (-[@{bond}:90,,,,Orange]@{hydrogen}\textcolor{Orange}{H})
                -[:-30]
                (=[:-90]O)
                -[:30]R
              }
            }{pK$_\text{a} \approx 20$}
            \qquad \qquad
            \arrow{<->>[][$\text{K} = 10^{(15-20)}$]}
            \qquad
            \chemname{
              \raisebox{-0.5cm}{
                \chemfig{
                  [:0]HO
                  -[:0,,,,Orange]\textcolor{Orange}{H}
                }
              }
            }{pK$_\text{a} \approx 15$}
            \raisebox{-0.5cm}{
              \+
              \chemfig{
                [:0]R
                -[:30]\ncharge{90}{}
                -[:-30]
                (=[:-90]O)
                -[:30]R
              }
            }
          }
          \chemmovearrow{
            \reactionfigarrow{ho}{45}{hydrogen}{180}{Red}
            \reactionfigarrow{bond}{15}{carbon}{10}{Red}
          }
        }{Enolate formation}{fig:ph5a}{0.9}

        \item \textbf{C-protonation of an enolate}

        \reactionfig{
          \scheme{
            \chemname{
              \raisebox{-0.5cm}{
                \chemfig{
                  [:0]@{oxygen1}HO
                  -[@{bond1}:0,,,,Orange]@{hydrogen}\textcolor{Orange}{H}
                }
              }
            }{pK$_\text{a} \approx 15$}
            \raisebox{-0.5cm}{
              \+
            }
            \chemfig{
              [:0]@{oxygen2}\ncharge{120}{O}
              -[@{bond2}:-90]
              =[@{bond3}:-30]@{carbon}
            }
            \qquad \qquad
            \arrow{<->>[][$\text{K} = 10^{(20-15)}$]}
            \qquad
            \chemname{
              \chemfig{
                [:0]O
                =[:-90]
                -[:-30]
                -[:-90,,,,Orange]\textcolor{Orange}{H}
              }
            }{pK$_\text{a} \approx 20$}
            \raisebox{-1cm}{
              \+
              \chemfig{
                [:0]\ce{HO-}
              }
            }
          }
          \chemmovearrow{
            \reactionfigarrow{bond1}{80}{oxygen1}{100}{Red}
            \reactionfigarrow{oxygen2}{45}{bond2}{15}{Red}
            \reactionfigarrow{bond3}{60}{carbon}{15}{Red}
            \reactionfigarrow{carbon}{-135}{hydrogen}{0}{Red}
          }
        }{C-protonation of an enolate}{fig:ph6a}{0.9}

        \item \textbf{Enol formation} 

        \reactionfig{
          \scheme{
            \chemname{
              \raisebox{-1.5cm}{
                \chemfig{
                  [:0]@{oxygen1}HO
                  -[@{bond}:0,,,,Orange]@{hydrogen}\textcolor{Orange}{H}
                }
              }
            }{pK$_\text{a} \approx 15$}
            \raisebox{-1.5cm}{
              \+
              \chemfig{
                [:0]
                (-[:90]@{oxygen2}\ncharge{135}{O})
                =[:-30]
              }
            }
            \qquad \qquad
            \arrow{<<->[][$\text{K} = 10^{(<0)}$]}
            \qquad
            \chemname{
              \raisebox{-1cm}{
                \chemfig{
                  [:0]\ncharge{45}{O}
                  (-[:150,,,,Orange]\textcolor{Orange}{H})
                  -[:-90]
                  =[:-30]
                }
              }
            }{pK$_\text{a} < 15$}
            \raisebox{-1.5cm}{
              \+
              \chemfig{
                [:0]\ce{HO-}
              }
            }
          }
          \chemmovearrow{
            \reactionfigarrow{oxygen2}{180}{hydrogen}{15}{Red}
            \reactionfigarrow{bond}{-45}{oxygen1}{-80}{Red}
          }
        }{Enol formation}{fig:ph7a}{0.9}

      \end{enumerate}

      Furthermore, pK$_\text{a}$ values of related structures can be predicted from the reference values in \autoref{tab:pka_table}, as demonstrated in the following examples.

      \begin{enumerate}

        \item \textbf{Amine like α to carbonyl} 

        \reactionfig{
          \scheme{
            \chemname{
              \raisebox{-1.5cm}{
                \chemfig{
                  [:0]\ce{H3C}
                  -[:-30]N
                  (-[:-90,,,,Red]\textcolor{Red}{H})
                  -[:30]
                  (=[:90]O)
                  -[:-30]\ce{CH3}
                }
              }
            }{$20 >$ pK$_\text{a} > 5$}
          }
        }{Amine like α to carbonyl}{fig:ph8a}{0.8}

        The above molecule has a higher pK$_\text{a}$ than a carboxylic acid (pK$_\text{a} = 5$), since nitrogen is less electronegative than oxygen, but a lower pK$_\text{a}$ than a proton α to a carbonyl group (pK$_\text{a} = 20$), since nitrogen is more electronegative than carbon.

        \item \textbf{α to carbonyl and allylic} 

        \reactionfig{
          \scheme{
            \chemname{
              \raisebox{-1.5cm}{
                \chemfig{
                  [:0]
                  -[:30]
                  (-[:90])
                  =[:-30]
                  -[:30]
                  (-[:90,,,,Red]\textcolor{Red}{H})
                  -[:-30]
                  (=[:-90]O)
                  -[:30]\ce{CH3}
                }
              }
            }{$20 >$ pK$_\text{a} > 10$}
          }
        }{α to carbonyl group and allylic}{fig:ph9a}{0.8}

        Here, The above molecule has a higher pK$_\text{a}$ than a proton α to two carbonyl groups (pK$_\text{a} = 10$), but a lower pK$_\text{a}$ than one α to a single carbonyl group (pK$_\text{a} = 20$), since it is additionally allylic.

      \end{enumerate}

    \end{core}

  \subsection{Allylic Systems}

    \begin{core}{Allylic Carbons}

      \textbf{Allylic carbons} are sp$^3$ hybridized carbons directly attached to a C=C double bond, and an \textbf{allylic hydrogen} is one bonded to such a carbon.

      \reactionfig{
        \scheme{
          \chemfig{
            [:0]
            -[:-30]\textcolor{Red}{C}
            (-[:-90]\textcolor{Blue}{H})
            (-[:150]\textcolor{Blue}{H})
            (-[:-150]\textcolor{Blue}{H})
            -[:30]
            =[:-30]
            -[:30]\textcolor{Red}{C}
            (-[:-30]\textcolor{Blue}{H})
            (-[:90]\textcolor{Blue}{H})
            (-[:30]\textcolor{Blue}{H})
          }
        }
      }{Allylic carbons (\textit{red}) and allylic hydrogens (\textit{blue})}{fig:ph10a}{0.8}

    \end{core}

    \begin{core}{Allylic System Resonance}

      An allylic hydrogen exhibits greater acidity than an both alkane and alkyne hydrogens, as resonance delocalizes the negative charge across the resulting allylic anion, thereby rendering the anion a comparatively weaker base.

      \vspace{-1cm}
      \doublereactionfig{
        \scheme{
          \chemname{
            \raisebox{-2cm}{
              \chemfig{
                [:0]
                -[:30]
                -[:-30]@{carbon1}
                -[@{bond1}:30]@{hydrogen1}H
              }
            }
          }{pK$_\text{a} \approx 50$}
          \raisebox{-2cm}{
            \+
            \chemfig[atom sep=2em]{
              [:0]@{nitrogen}\ncharge{90}{N}
              (
              -[:-30]
              (-[:-70])
              -[:10]
              )
              (
              -[:-150]
              (-[:-110])
              -[:-190]
              )
            }
          }
          \qquad \qquad
          \arrow{<<->[][$\text{K}_\text{eq} = 10^{(35-50)}$]}
          \qquad
          \chemname{
            \raisebox{-2cm}{
              \chemfig[atom sep=2em]{
                [:0]N
                (-[:90]H)
                (
                -[:-30]
                (-[:-70])
                -[:10]
                )
                (
                -[:-150]
                (-[:-110])
                -[:-190]
                )
              }
            }
          }{pK$_\text{a} \approx 35$}
          \raisebox{-2cm}{
            \+
            \chemfig{
              [:0]
              -[:30]
              -[:-30]\ncharge{45}{}
            }
          }
        }
        \chemmovearrow{
          \reactionfigarrow{bond1}{-45}{carbon1}{-90}{Red}
          \reactionfigarrow{nitrogen}{135}{hydrogen1}{45}{Red}
        }
      }{
        \scheme{
          \chemname{
            \raisebox{-2cm}{
              \chemfig{
                [:0]
                =[:30]
                -[:-30]@{carbon1}
                -[@{bond1}:30]@{hydrogen1}H
              }
            }
          }{pK$_\text{a} \approx 40$}
          \raisebox{-2cm}{
            \+
            \chemfig[atom sep=2em]{
              [:0]@{nitrogen}\ncharge{90}{N}
              (
              -[:-30]
              (-[:-70])
              -[:10]
              )
              (
              -[:-150]
              (-[:-110])
              -[:-190]
              )
            }
          }
          \qquad \qquad
          \arrow{<<->[][$\text{K}_\text{eq} = 10^{(35-40)}$]}
          \qquad
          \chemname{
            \raisebox{-2cm}{
              \chemfig[atom sep=2em]{
                [:0]N
                (-[:90]H)
                (
                -[:-30]
                (-[:-70])
                -[:10]
                )
                (
                -[:-150]
                (-[:-110])
                -[:-190]
                )
              }
            }
          }{pK$_\text{a} \approx 35$}
          \raisebox{-2cm}{
            \+
            \chemfig{
              [:0]@{carbon2}
              =[@{dbond}:30]
              -[@{bond2}:-30]@{negative}\ncharge{-90}{}
            }
          }
          \arrow{<->}
          \chemfig{
            [:0]\ncharge{-90}{}
            -[:30]
            =[:-30]
          }
        }
        \chemmovearrow{
          \reactionfigarrow{bond1}{-45}{carbon1}{-90}{Red}
          \reactionfigarrow{nitrogen}{135}{hydrogen1}{45}{Red}
          \reactionfigarrow{negative}{30}{bond2}{60}{Red}
          \reactionfigarrow{dbond}{120}{carbon2}{150}{Red}
        }
      }{K$_\text{eq}$s of alkane (\textit{top}) vs allylic (\textit{bottom}) hydrogens; \textit{the K$_\text{eq}$ of the allylic scenario is much more forward favoring because it is abel to better stablize the negative anion state charge}}{fig:ph11a}{0.72}
      \vspace{-0.5cm}

      Just as negative charge (\textit{filled octets}) is stabilized by resonance within allylic structures, radicals and positive charge (\textit{unfilled octets}) are stabilized through this same resonance mechanism.

      \triplereactionfig{
        \scheme{
          \chemfig{
            [:0]@{carbon2}
            =[@{dbond}:30]
            -[@{bond2}:-30]@{negative}\ncharge{-90}{}
          }
          \arrow{<->}
          \chemfig{
            [:0]\ncharge{-90}{}
            -[:30]
            =[:-30]
          }
        }
        \chemmovearrow{
          \reactionfigarrow{negative}{30}{bond2}{60}{Red}
          \reactionfigarrow{dbond}{120}{carbon2}{150}{Red}
        }
      }{
        \scheme{
          \chemfig{
            [:0]@{carbon2}
            =[@{dbond}:30]
            -[@{bond2}:-30]@{negative}\charge{-90:4pt=\LARGE$\cdot$}{}
          }
          \arrow{<->}
          \chemfig{
            [:0]\charge{-90:4pt=\LARGE$\cdot$}{}
            -[:30]
            =[:-30]
          }
        }
        \chemmovearrow{
          \reactionfigarrow{negative}{30}{bond2}{60}{Green}
          \reactionfigarrow{dbond}{120}{carbon2}{150}{Green}
          \reactionfigarrow{dbond}{90}{bond2}{60}{Green}
        }
      }{
        \scheme{
          \chemfig{
            [:0]
            =[@{dbond}:30]
            -[@{bond2}:-30]\pcharge{-90}{}
          }
          \arrow{<->}
          \chemfig{
            [:0]\pcharge{-90}{}
            -[:30]
            =[:-30]
          }
        }
        \chemmovearrow{
          \reactionfigarrow{dbond}{120}{bond2}{60}{Red}
        }
      }{Filled octets (\textit{top}), radicals (\textit{middle}), and unfilled octets (\textit{bottom})}{fig:ph12a}{0.72}
      \vspace{-0.5cm}

      Critically, \textbf{any lone pair capable of participating in resonance must reside in a p orbital}.

    \end{core}

    \begin{core}{Allylic System Bond Dissociation Energy}

      As a consequence of resonance-induced stabilization, the bond dissociation energies of distinct hydrogens within an alkene follow a predictable trend.

      \vspace{-0.15cm}
      \triplereactionfig{
        \scheme{
          \text{(1$^\circ$ CH)}
          \qquad
          \chemfig[atom sep=2em]{
            [:0]
            =[:30]
            -[:-30]
            -[:30]
            (-[:90])
            -[:-30]@{carbon}
            -[@{bond}:30,,,,Orange]@{hydrogen}\textcolor{Orange}{H}
          }
          \qquad \qquad
          \arrow{->[$+101$]}
          \qquad
          \chemfig[atom sep=2em]{
            [:0]
            =[:30]
            -[:-30]
            -[:30]
            (-[:90])
            -[:-30]\textcolor{Orange}{\charge{-90:4pt=\LARGE$\cdot$}{}}
          }
          \+
          \chemfig{
            [:0]\textcolor{Orange}{H\cdot}
          }
        }
        \chemmovearrow{
          \reactionfigarrow{bond}{-45}{hydrogen}{-30}{Green}
          \reactionfigarrow{bond}{90}{carbon}{-150}{Green}
        }
      }{
        \scheme{
          \text{(3$^\circ$ CH)}
          \qquad
          \chemfig[atom sep=2em]{
            [:0]
            =[:30]
            -[:-30]
            -[:30]@{carbon}
            (-[:90])
            (-[@{bond}:30,,,,Orange]@{hydrogen}\textcolor{Orange}{H})
            -[:-30]
          }
          \qquad \qquad
          \arrow{->[$+96.5$]}
          \qquad
          \chemfig[atom sep=2em]{
            [:0]
            =[:30]
            -[:-30]
            -[:30]\textcolor{Orange}{\charge{-90:4pt=\LARGE$\cdot$}{}}
            (-[:90])
            -[:-30]
          }
          \+
          \chemfig{
            [:0]\textcolor{Orange}{H\cdot}
          }
        }
        \chemmovearrow{
          \reactionfigarrow{bond}{100}{hydrogen}{80}{Green}
          \reactionfigarrow{bond}{-60}{carbon}{-90}{Green}
        }
      }{
        \scheme{
          \text{(2$^\circ$ CH \& Allylic)}
          \qquad
          \chemfig[atom sep=2em]{
            [:0]
            =[:30]
            -[:-30]@{carbon}
            (-[@{bond}:-90,,,,Orange]@{hydrogen}\textcolor{Orange}{H})
            -[:30]
            (-[:90])
            -[:-30]
          }
          \qquad \qquad
          \arrow{->[$+87$]}
          \qquad
          \chemfig[atom sep=2em]{
            [:0]
            =[:30]
            -[:-30]\textcolor{Orange}{\charge{90:4pt=\LARGE$\cdot$}{}}
            -[:30]
            (-[:90])
            -[:-30]
          }
          \+
          \chemfig{
            [:0]\textcolor{Orange}{H\cdot}
          }
        }
        \chemmovearrow{
          \reactionfigarrow{bond}{10}{hydrogen}{-10}{Green}
          \reactionfigarrow{bond}{-170}{carbon}{120}{Green}
        }
      }{Primary (\textit{top}), tertiary (\textit{middle}), and secondary \& allylic (\textit{bottom}) hydrogen bond dissociation energies within the same alkene; \textit{secondary allylic hydrogen bond takes least energy to break due to resonance-induced stabilization}}{fig:ph13a}{0.72}

      Above, the 3$^\circ$ carbon radical is lower in energy than the 1$^\circ$ carbon radical as a result of hyperconjugation between the radical carbon and neighboring $\sigma$ bonds. The allylic radical, by contrast, receives comparatively less hyperconjugative stabilization as a 2$^\circ$ carbon, yet is stabilized through resonance, causing its corresponding bond dissociation to require the least energy of the three.

      \pgfig{
        \draw[->] (0,0) -- (0,3.2);
        \node[] at (-0.4,1.6) {$E$};

        % --- sp3 panel ---
        \begin{scope}[xshift=1cm]
          \draw[black] (0.2,2.5) -- (1.2,2.5);
          \draw[black,->] (0.7,0.5) -- (0.7,2.4);
          \node at (0.7,0.3) {1$^\circ$ rad};
          \node at (0.3,1.45) {101};
        \end{scope}

        % --- sp2 panel ---
        \begin{scope}[xshift=3cm]
          \draw[black] (0.2,2.0) -- (1.2,2.0);
          \draw[black,->] (0.7,0.5) -- (0.7,1.9);
          \node at (0.7,0.3) {3$^\circ$ rad};
          \node at (0.3,1.2) {96.5};
        \end{scope}

        % --- sp panel ---
        \begin{scope}[xshift=5cm]
          \draw[black] (0.2,1.5) -- (1.2,1.5);
          \draw[black,->] (0.7,0.5) -- (0.7,1.4);
          \node at (0.7,0.3) {2$^\circ$ allylic rad};
          \node at (0.3,0.95) {87};
        \end{scope}
      }{Energy diagram of \autoref{fig:ph13a}}{fig:ph14a}{0.95}

      This differential stabilization gives rise to chemical consequences that play a substantial role in practice, as illustrated by the reaction below.

      \doublereactionfig{
        \scheme{
          \chemfig{
            [:0]
            -[:30]
            -[:-30]
            -[@{bond}:30]@{ots}OTs
          }
          \+
          \chemfig{
            \ce{MeOH}
          }
          \arrow{->[MeOH][Δ]}
          \text{No Reaction; \textit{\textcolor{Red}{primary carbocation generation}}}
        }
        \chemmovearrow{
          \reactionfigarrow{bond}{-80}{ots}{-50}{dash pattern= on 2pt off 2pt,Red}
        }
      }{
        \scheme{
          \chemfig{
            [:0]
            =[:30]
            -[:-30]
            -[@{bond}:30]@{ots}OTs
          }
          \+
          \chemfig{
            \ce{MeOH}
          }
          \arrow(aa--bb){->[MeOH][Δ]}
          \chemfig{
            [:0]
            =[:30]
            -[:-30]
            -[:30]OMe
          }
          \+
          \ce{TsO-}
          \+
          \ce{MeOH2+}
          \arrow(@aa--cc)[-90,1.2]
          \chemfig[atom sep=2em]{
            [:0]
            =[:30]
            -[:-30]@{positive}\pcharge{-90}{}
          }
          \+
          \ce{TsO-}
          \arrow(@cc--dd){->[
            \chemfig{
              @{meoh}\ce{MeOH}
            }
          ][]}[0]
          \chemfig[atom sep=2em]{
            [:0]
            =[:30]
            -[:-30]
            -[:30]@{oxygen}\pcharge{-90}{O}
            (-[@{bond2}:90]@{hydrogen}H)
            -[:-30]Me
          }
          \+
          \chemfig{
            @{meoh2}\ce{MeOH}
          }
          \arrow(@dd--@bb)
        }
        \chemmovearrow{
          \reactionfigarrow{bond}{-80}{ots}{-50}{Red}
          \reactionfigarrow{meoh}{150}{positive}{60}{Red}
          \reactionfigarrow{meoh2}{100}{hydrogen}{30}{Red}
          \reactionfigarrow{bond2}{15}{oxygen}{0}{Red}
        }
      }{Non-allylic (\textit{top}) vs allylic (\textit{bottom}) leaving group primary carbocation formation viability; \textit{a primary carbocation can form in the bottom reaction because the cation is additionally allylic}; \textit{Ts stands for a tosyl group, a common chemical group used in organic sythesis}}{fig:ph15a}{0.9}

    \end{core}

    \begin{core}{Allylic System Molecular Orbital Structure}

      The molecular orbitals \textbf{relevant to allylic structures are the p orbitals}, whose energy levels can be compared in the context of sp- and sp$^2$-hybridized carbons, as shown below. Notably, the energy level of sp$^2$ orbitals lies higher than that of sp orbitals.

      \pgfig{
        \node at (0.6,2.5) {sp};
        \draw[->] (0,0) -- (0,2.2);
        \node at (-0.4,1.1) {$E$};

        \draw (0.2,1.7) -- (0.5,1.7);
        \draw (0.6,1.7) -- (0.9,1.7);
        \node at (1.2,1.7) {p};

        \draw (0.2,0.6) -- (0.5,0.6);
        \draw (0.6,0.6) -- (0.9,0.6);
        \node at (1.2,0.6) {sp};
        \draw[dashed] (1.6,0.6) -- (6.3,0.6);

        \begin{scope}[xshift=5cm]
          \node at (0.6,2.5) {sp$^2$};
          \draw[->] (0,0) -- (0,2.2);
          \node at (-0.4,1.1) {$E$};

          \draw (0.2,1.7) -- (0.5,1.7);
          \node at (0.8,1.7) {p};

          \draw (0.2,1.0) -- (0.5,1.0);
          \draw (0.6,1.0) -- (0.9,1.0);
          \draw (1.0,1.0) -- (1.3,1.0);
          \node at (1.6,1.0) {sp$^2$};
        \end{scope}
      }{Orbital energy level diagrams for sp and sp$^2$ hybridization}{fig:ph16a}{0.95}

      Superimposed on a skeletal diagram backbone, this orbital structure appears as follows.

      \pgfig{
        \node (mol) {
          \chemfig{
            [:0]
            (<:[:150]H)
            (<[:-120]H)
            -[:15]
            -[:-15]\charge{-20:8pt=$\scriptstyle\ominus$}{}\charge{75:20pt=\:}{}
            (<:[:30]H)
            (<[:-60]H)
          }
        };

        % middle
        \draw[blue, rotate=0] (0,1.15) ellipse (0.2cm and 0.6cm);
        \draw[blue, rotate=0] (0,-0.25) ellipse (0.2cm and 0.6cm);

        % right
        \draw[blue, rotate=-15] (0.95,1.15) ellipse (0.2cm and 0.6cm);
        \draw[blue, rotate=-15] (0.95,-0.25) ellipse (0.2cm and 0.6cm);

        % left
        \draw[blue, rotate=15] (-0.95,1.15) ellipse (0.2cm and 0.6cm);
        \draw[blue, rotate=15] (-0.95,-0.25) ellipse (0.2cm and 0.6cm);
      }{σ bond orbitals overlap along internuclear axis (handshake) whereas π bond orbitals overlap above and below internuclear axes. All cabons are sp2 hybridized}{fig:ph17a}{0.9}

      Considered as a whole, the allylic p-orbital system produces an energy diagram with a distinctive three-level pattern, progressing from bonding to non-bonding to anti-bonding.

      \myfig{0.55}{conj_sys_figures/allylic_orbital}{Energy diagram of allylic p-orbital system}{fig:ph18}

      Consequently, the HOMO and LUMO assignments differ among the allylic anion, radical, and cation.

      \pgfig{
        % ================= Energy axis =================
        \draw[thick,-{Latex}] (-1.5,-0.5) -- (-1.5,4.5) node[above]{}; 
        \node[rotate=0] at (-2.1,2) {$E$};
        % ================= Column 1: Allyl anion (4 e-) =================
        \begin{scope}[shift={(0,0)}]
          % skeleton placeholder
          \node at (0.6,4.3) {
            \chemfig[atom sep=2em]{
              [:0]
              =[:30]
              -[:-30]\ncharge{-90}{}
            }
          };
          % pi levels (halved width: 0 to 0.6)
          \draw[thick] (0,3) -- (0.6,3); \node[right] at (0.6,3) {π3};
          \node[orange] at (1.6,3) {LUMO};
          \draw[thick] (0,1.8) -- (0.6,1.8); \node[right] at (0.6,1.8) {π2};
          \node[orange] at (1.6,1.8) {HOMO};
          \draw[thick] (0,0.6) -- (0.6,0.6); \node[right] at (0.6,0.6) {π1};
          % electrons
          \draw[Blue,-{Latex}] (0.15,1.5) -- (0.15,2.1);
          \draw[Blue,-{Latex}] (0.45,2.1) -- (0.45,1.5);
          \draw[Blue,-{Latex}] (0.15,0.3) -- (0.15,0.9);
          \draw[Blue,-{Latex}] (0.45,0.9) -- (0.45,0.3);
        \end{scope}
        % ================= Column 2: Allyl radical (3 e-) =================
        \begin{scope}[shift={(4,0)}]
          \node at (0.6,4.3) {
            \chemfig[atom sep=2em]{
              [:0]
              =[:30]
              -[:-30]\charge{-90:4pt=$\cdot$}{}
            }
          };
          \draw[thick] (0,3) -- (0.6,3); \node[right] at (0.6,3) {π3};
          \node[orange] at (1.6,3) {LUMO};
          \draw[thick] (0,1.8) -- (0.6,1.8); \node[right] at (0.6,1.8) {π2};
          \node[orange] at (1.6,1.8) {HOMO};
          \draw[thick] (0,0.6) -- (0.6,0.6); \node[right] at (0.6,0.6) {π1};
          \draw[Blue,-{Latex}] (0.3,1.5) -- (0.3,2.1);
          \draw[Blue,-{Latex}] (0.15,0.3) -- (0.15,0.9);
          \draw[Blue,-{Latex}] (0.45,0.9) -- (0.45,0.3);
        \end{scope}
        % ================= Column 3: Allyl cation (2 e-) =================
        \begin{scope}[shift={(8,0)}]
          \node at (0.6,4.3) {
            \chemfig[atom sep=2em]{
              [:0]
              =[:30]
              -[:-30]\pcharge{-90}{}
            }
          };
          \draw[thick] (0,3) -- (0.6,3); \node[right] at (0.6,3) {π3};
          \draw[thick] (0,1.8) -- (0.6,1.8); \node[right] at (0.6,1.8) {π2};
          \node[orange] at (1.6,1.8) {LUMO};
          \draw[thick] (0,0.6) -- (0.6,0.6); \node[right] at (0.6,0.6) {π1};
          \node[orange] at (1.6,0.6) {HOMO};
          \draw[Blue,-{Latex}] (0.15,0.3) -- (0.15,0.9);
          \draw[Blue,-{Latex}] (0.45,0.9) -- (0.45,0.3);
        \end{scope}
      }{Molecular orbital diagrams for the allyl anion, radical, and cation, showing the filling of π1, π2, and π3 with 4, 3, and 2 π \ce{e^-}s respectively, along with HOMO LUMO assignments}{fig:ph20a}{0.9}

      This distinction is particularly relevant in the context of \textbf{frontier molecular orbital theory}, which holds that the HOMO of a nucleophile interacts with the LUMO of an electrophile (\textit{or vice versa}) as explored in the context of diel alder reaction in UNCOMMENT AT END %\autoref{core:da_orbitals}.

    \end{core}

  \subsection{Conjugated π Systems}

    \begin{core}{Conjugated π System Stability}

      A \textbf{conjugated π system} refers to a network of adjacent allylic structures that permits resonance delocalization across a greater number of carbons. This effect can be observed by examining the hydrogenation enthalpies of the following alkenes.

      \begin{figure}[H]
        \centering

        \begin{minipage}{\textwidth}
          \centering
          \scalebox{0.9}{
            \scheme{
              \chemfig[atom sep=2em]{
                [:0]
                -[:30]
                -[:-30]
                -[:30]
                =[:-30]
                -[:30]
              }
              \arrow{->[\ce{H2}]}
              \chemfig[atom sep=2em]{
                [:0]
                -[:30]
                -[:-30]
                -[:30]
                -[:-30]
                -[:30]
              }
              \qquad
              \chemfig{
                \text{$ΔH_\text{rxn}=-27.7$}
              }
            }
          }
        \end{minipage}

        \vspace{2.5em}

        \begin{minipage}{\textwidth}
          \centering
          \scalebox{0.9}{
            \scheme{
              \chemfig[atom sep=2em]{
                [:0]
                -[:30]
                -[:-30]
                =[:30]
                -[:-30]
                -[:30]
              }
              \arrow{->[\ce{H2}]}
              \chemfig[atom sep=2em]{
                [:0]
                -[:30]
                -[:-30]
                -[:30]
                -[:-30]
                -[:30]
              }
              \qquad
              \chemfig{
                \text{$ΔH_\text{rxn}=-29.9$}
              }
            }
          }
        \end{minipage}

        \vspace{2.5em}

        \begin{minipage}{\textwidth}
          \centering
          \scalebox{0.9}{
            \scheme{
              \chemfig[atom sep=2em]{
                [:0]
                -[:30]
                -[:-30]
                -[:30]
                -[:-30]
                =[:30]
              }
              \arrow{->[\ce{H2}]}
              \chemfig[atom sep=2em]{
                [:0]
                -[:30]
                -[:-30]
                -[:30]
                -[:-30]
                -[:30]
              }
              \qquad
              \chemfig{
                \text{$ΔH_\text{rxn}=-28.2$}
              }
            }
          }
        \end{minipage}

        \vspace{2.5em}

        \begin{minipage}{\textwidth}
          \centering
          \scalebox{0.9}{
            \scheme{
              \chemfig[atom sep=2em]{
                [:0]
                =[:30]
                -[:-30]
                -[:30]
                =[:-30]
                -[:30]
              }
              \arrow{->[\ce{2H2}]}
              \chemfig[atom sep=2em]{
                [:0]
                -[:30]
                -[:-30]
                -[:30]
                -[:-30]
                -[:30]
              }
              \qquad
              \chemfig{
                \shortstack{
                  Predict: $ΔH_\text{1}+ΔH_\text{3}=-57.6$\\
                  Actual: $ΔH_\text{rxn}=-57.6\,\checkmark$
                }
              }
            }
          }
        \end{minipage}

        \vspace{2.5em}

        \begin{minipage}{\textwidth}
          \centering
          \scalebox{0.9}{
            \scheme{
              \chemfig[atom sep=2em]{
                [:0]
                =[:30]
                -[:-30]
                =[:30]
                -[:-30]
                -[:30]
              }
              \arrow{->[\ce{2H2}]}
              \chemfig[atom sep=2em]{
                [:0]
                -[:30]
                -[:-30]
                -[:30]
                -[:-30]
                -[:30]
              }
              \qquad
              \chemfig{
                \shortstack{
                  Predict: $ΔH_\text{2}+ΔH_\text{3}=-58.1$\\
                  Actual: \textcolor{Red}{$ΔH_\text{rxn}=-52.8$}
                }
              }
            }
          }
        \end{minipage}

        \caption{Enthalpies of alkene and alkadiene hydrogenations; \textit{the actual enthalpy of a conjugated π system alkadiene is lower than the predicted sum of individual bond hydrogenation enthalpies}}
        \label{fig:ph21a}
      \end{figure}

      Above, compound 5 exhibits greater stabilization than predicted by the simple additive sum of its constituent enthalpies, owing to the enhanced ability of conjugated π systems to stabilize excess or deficient charge.

      \pgfig{
        % ================= Panel 4: Predicted matches Actual =================
        \begin{scope}[shift={(0,0)}]
          \node at (-0.8,4.6) {\large (4)};
          \draw[thick,-{Latex}] (-0.8,0) -- (-0.8,4.3);
          \node at (-1.4,2.2) {$E$};

          % starting material level (diene)
          \draw[thick] (0,3.6) -- (1,3.6);
          \node at (0.5,4.0) {
            \chemfig[atom sep=1em]{
              [:0]
              =[:30]
              -[:-30]
              -[:30]
              =[:-30]
              -[:30]
            }
          };

          % predicted arrow (black)
          \draw[thick,-{Latex}] (0.5,3.5) -- (0.5,0.6);
          \node[black, align=right, font=\small] at (0.0,2.1) {-57.6\\ Pred};

          % predicted product level
          \draw[thick] (0,0.5) -- (1,0.5);
          \node at (0.5,0.15) {
            \chemfig[atom sep=1em]{
              [:0]
              -[:30]
              -[:-30]
              -[:30]
              -[:-30]
              -[:30]
            }
          };

          % actual (yellow), same energy as predicted
          \draw[yellow!80!black, thick] (2,3.6) -- (3,3.6);
          \node[yellow!70!black] at (2.5,4.0) {
            \chemfig[atom sep=1em]{
              [:0]
              =[:30]
              -[:-30]
              -[:30]
              =[:-30]
              -[:30]
            }
          };
          \draw[yellow!80!black, thick,-{Latex}] (2.5,3.5) -- (2.5,0.6);
          \node[yellow!70!black, align=right, font=\small] at (1.9,2.1) {-57.6\\ Actual};
          \draw[yellow!80!black, thick] (2,0.5) -- (3,0.5);
          \node[yellow!70!black] at (2.5,0.15) {
            \chemfig[atom sep=1em]{
              [:0]
              -[:30]
              -[:-30]
              -[:30]
              -[:-30]
              -[:30]
            }
          };
        \end{scope}

        % ================= Panel 5: Actual is stabilized relative to Predicted =================
        \begin{scope}[shift={(6,0)}]
          \node at (-0.8,4.6) {\large (5)};
          \draw[thick,-{Latex}] (-0.8,0) -- (-0.8,4.3);
          \node at (-1.4,2.2) {$E$};

          % starting material level (diene)
          \draw[thick] (0,3.6) -- (1,3.6);
          \node at (0.5,4.0) {
            \chemfig[atom sep=1em]{
              [:0]
              =[:30]
              -[:-30]
              =[:30]
              -[:-30]
              -[:30]
            }
          };

          % predicted arrow (black)
          \draw[thick,-{Latex}] (0.5,3.5) -- (0.5,0.6);
          \node[black, align=right, font=\small] at (0.0,2.1) {-58.1\\ Pred};

          % predicted product level (black, lower)
          \draw[thick] (0,0.5) -- (1,0.5);
          \node at (0.5,0.15) {
            \chemfig[atom sep=1em]{
              [:0]
              -[:30]
              -[:-30]
              -[:30]
              -[:-30]
              -[:30]
            }
          };

          % actual (yellow), same energy as predicted
          \draw[yellow!80!black, thick] (2,2.5) -- (3,2.5);
          \node[yellow!70!black] at (2.5,2.9) {
            \chemfig[atom sep=1em]{
              [:0]
              =[:30]
              -[:-30]
              =[:30]
              -[:-30]
              -[:30]
            }
          };
          \draw[yellow!80!black, thick,-{Latex}] (2.5,2.4) -- (2.5,0.6);
          \node[yellow!70!black, align=right, font=\small] at (1.9,1.65) {-52.8\\ Actual};
          \draw[yellow!80!black, thick] (2,0.5) -- (3,0.5);
          \node[yellow!70!black] at (2.5,0.15) {
            \chemfig[atom sep=1em]{
              [:0]
              -[:30]
              -[:-30]
              -[:30]
              -[:-30]
              -[:30]
            }
          };

          % stabilization double arrow between black and yellow product levels
          \draw[orange, thick,{Latex}-{Latex}] (1.2,3.6) -- (1.2,2.5);


        \end{scope}
      }{Comparison of predicted versus experimentally measured heats of hydrogenation for a non-conjugated diene (\textit{left}) and a conjugated diene (\textit{right}); \textit{the conjugated system's smaller actual value reflects a 5 kcal/mol stabilization arising from the four contiguous p-orbitals of the conjugated π system}}{fig:ph22a}{1}
      \vspace{-0.4cm}

      A closer examination of the orbital energy diagrams for these conjugated π systems, following the same approach applied previously to the three-orbital π system (\textit{\autoref{fig:ph18}}), reveals that the conjugated system's orbital energy levels span an overall larger range while exhibiting smaller stepwise increases in energy between successive levels. Consequently, the HOMO--LUMO gap is comparatively smaller than that of the non-conjugated alkene π system.

      \myfig{0.65}{conj_sys_figures/conj_edia}{Energy diagrams of non-conjugated (\textit{left}) and conjugated (\textit{right}) π systems in alkadienes from \autoref{fig:ph22a}; \textit{source of stability in a conjugated π system are the lowest π MO \ce{e^-}s}}{fig:ph23a}
      \vspace{-0.4cm}

      The governing principles for conjugated π system orbital energy digrams are as follows.

      \begin{enumerate}

        \item As conjugation increases (i.e., more p orbitals arranged in succession), the energy gap between π MOs decreases.

        \item As conjugation increases, the lowest-energy π MO decreases further in energy, while the highest-energy π MO increases further in energy.

      \end{enumerate}

      It is important to recognize that every high-energy intermediate is stabilized to the greatest extent possible, and that resonance stabilization arising from conjugated π systems supersedes stabilization originating from hyperconjugation.
      \[
        \text{resonance} > 3^\circ > 2^\circ > 2^\circ
      \]

      This principle carries particularly significant consequences for the reactions covered within \href{https://ibahng.github.io/chemistry/alkene_alkyne/alkene_alkyne.pdf}{Alkenes \& Alkynes}, as many electrophilic additions proceed in either a Markovnikov or anti-Markovnikov manner, a distinction that depends on the relative stability of the intermediate state. Representative examples include alkene hydrohalogenation (\textit{\autoref{alk-core:hydrohalogenation}}), which proceeds in a Markovnikov manner, and hydroboration-oxidation (\textit{\autoref{alk-core:alkene_hydro_oxi}}), which proceeds in an anti-Markovnikov manner.

    \end{core}

    Applying our understanding of conjugated π system stability, we now consider two illustrative examples.

    \begin{core}[label=core:chlorination_example]{Chlorination of Conjugated Cycloalkene}

      Our objective is to predict all products of the following reaction, including any relevant stereoisomers.

      \reactionfig{
        \scheme{
          \chemfig[atomsep=1.5em]{
            [:0]*5(=
            (-[:-70])
            -
            (-[:30])
            (-[:-30])
            -=-)
          }
          \+
          \chemfig{
            \ce{Cl2}
          }
          \arrow{->[\ce{H2O}][]}
          \chemfig{
            \text{\textit{predict all products including stereoisomers}}
          }
        }
      }{Cycloalkene chlorination reactants}{fig:ph24a}{0.85}

      First, we map out the chloronium intermediates.

      \reactionfig{
        \scheme{
          \chemfig[atomsep=1.5em]{
            [:0]*5(@{carbonleft}=[@{bond1}:]
            @{carbonright}
            (
            -[:-70]
            -[:-150,,,,White]@{cltop}Cl
            -[@{bond2}:-90]@{clbottom}Cl
            )
            -
            (-[:30])
            (-[:-30])
            -=-)
          }
          \arrow{->}
          \chemfig[atomsep=1.5em]{
            [:-72]*5(
            -
            (*3(<\pcharge{-90}{Cl}>-))
            -
            (<:[:-70])
            -
            (-[:30])
            (-[:-30])
            -=
            )
          }
          \raisebox{-1cm}{
            \+
          }
          \chemfig[atomsep=1.5em]{
            [:-72]*5(
            -
            (*3(<:\pcharge{-90}{Cl}>:-))
            -
            (<[:-70])
            -
            (-[:30])
            (-[:-30])
            -=
            )
          }
        }
        \chemmovearrow{
          \reactionfigarrow{bond2}{10}{clbottom}{-10}{Red}
          \reactionfigarrow{cltop}{110}{carbonleft}{-80}{Red}
          \reactionfigarrow{bond1}{-100}{carbonright}{-120}{Red}
          \reactionfigarrow{carbonright}{-100}{cltop}{70}{Red}
        }
      }{Top (\textit{left}) and bottom (\textit{back}) attacks leading to corresponding chloronium intermediates}{fig:ph25a}{0.85}

      Resonance is used to determine which carbon of the halonium intermediate is preferentially attacked by the weak nucleophile ROH.

      \reactionfig{
        \scheme{
          \chemfig[atomsep=1.5em]{
            [:-72]*5(
            -
            (*3(
            <@{cl1}\pcharge{-90}{Cl}
            >[@{bond1}:]
            -
            ))
            -
            (<:[:-70])
            -
            (-[:30])
            (-[:-30])
            -=
            )
          }
          \arrow{<->}
          \chemfig[atomsep=1.5em]{
            [:-72]*5(
            -
            (<[:-90]Cl)
            -\pcharge{110}{}
            (<:[:-70])
            -
            (-[:30])
            (-[:-30])
            -=
            )
          }
          \qquad
          \raisebox{-1cm}{
            \text{vs}
          }
          \qquad
          \chemfig[atomsep=1.5em]{
            [:-72]*5(
            -
            (*3(
            <[@{bond2}:]@{cl2}\pcharge{-90}{Cl}
            >
            -
            ))
            -
            (<:[:-70])
            -
            (-[:30])
            (-[:-30])
            -=
            )
          }
          \arrow{<->}
          \chemfig[atomsep=1.5em]{
            [:-72]*5(
            -[@{bond4}:]\pcharge{30}{}
            -
            (<[:-120]Cl)
            (<:[:-70])
            -
            (-[:30])
            (-[:-30])
            -=[@{bond3}:]
            )
          }
          \arrow{<->}
          \chemfig[atomsep=1.5em]{
            [:-72]*5(
            =
            -
            (<[:-120]Cl)
            (<:[:-70])
            -
            (-[:30])
            (-[:-30])
            -\pcharge{90}{}
            -
            )
          }
        }
        \chemmovearrow{
          \reactionfigarrow{bond1}{-30}{cl1}{-70}{Red}
          \reactionfigarrow{bond2}{-160}{cl2}{-110}{Red}
          \reactionfigarrow{bond3}{-45}{bond4}{-15}{Red}
        }
      }{Tertiary ``σ bond resonance" (\textit{left}) vs secondary and allylic (\textit{right})}{fig:ph26a}{0.7}

      Because the allylic resonance contributor is more stable, the electrophilic carbons are consequently located at two distinct positions within the system.

      \vspace{0.4cm}
      \reactionfig{
        \scheme{
          \chemfig[atomsep=1.5em]{
            [:-72]*5(
            -@{ecar2}\charge{-180:6pt=\textcolor{Blue}{δ+}}{}
            (*3(-\pcharge{-90}{}\charge{-180:6pt=\textcolor{Blue}{δ++}}{Cl}
            --))
            -
            (-[:-70])
            -
            (-[:30])
            (-[:-30])
            -@{ecar1}\charge{90:6pt=\textcolor{Blue}{δ+}}{}
            =
            )
          }
          \qquad \qquad
          \raisebox{-1cm}{
            \chemfig{
              @{ecar}\text{\textcolor{Red}{electrophilic carbons}}
            }
          }
        }
        \chemmovearrow{
          \reactionfigarrow{ecar}{135}{ecar1}{45}{Red}
          \reactionfigarrow{ecar}{-160}{ecar2}{30}{Red}
        }
      }{Electrophilic carbons on chloronium intermediate; \textit{the top resonance contributer is the Cl}}{fig:ph27a}{0.9}

      Consequently, three possible modes of attack exist, giving rise to three distinct stereoisomeric products.

      \begin{enumerate}

        \item \textbf{Tertiary electrophilic carbon attack} 

        \reactionfig{
          \scheme{
            \chemfig{
              @{water1}\ce{H2O}
            }
            \qquad \qquad
            \chemfig[atomsep=1.5em]{
              [:-72]*5(
              -@{carbon}
              (*3(
              <[@{bond}:]@{chlorine}\pcharge{-90}{Cl}
              >
              -
              ))
              -
              (<:[:-70])
              -
              (-[:30])
              (-[:-30])
              -=
              )
            }
            \arrow{->}
            \chemfig[atomsep=1.5em]{
              [:-72]*5(
              -
              (
              <:[:-120]@{oxygen}\pcharge{-20}{O}
              (-[:-170]H)
              (-[@{hbond}:-90]@{hydrogen}H)
              )
              -
              (<[:-120]Cl)
              (<:[:-70])
              -
              (-[:30])
              (-[:-30])
              -=[@{bond3}:]
              )
            }
            \arrow{->[][
              \chemfig{
                @{water2}\ce{H2O}
              }
            ]}
            \chemfig[atomsep=1.5em]{
              [:-72]*5(
              -
              (
              <:[:-120]HO
              )
              -
              (<[:-120]Cl)
              (<:[:-70])
              -
              (-[:30])
              (-[:-30])
              -=[@{bond3}:]
              )
            }
          }
          \chemmovearrow{
            \reactionfigarrow{water1}{-60}{carbon}{180}{Red}
            \reactionfigarrow{bond}{-170}{chlorine}{-110}{Red}
            \reactionfigarrow{hbond}{-120}{oxygen}{-150}{Red}
            \reactionfigarrow{water2}{-135}{hydrogen}{-10}{Red}
          }
        }{Tertiary electrophilic carbon attack}{fig:ph28a}{0.8}

        \item \textbf{Secondary electrophilic carbon attack from front} 

        \reactionfig{
          \scheme{
            \chemfig{
              @{water1}\ce{H2O}
            }
            \qquad \qquad
            \chemfig[atomsep=1.5em]{
              [:-72]*5(
              -[@{bond1}:]
              (*3(
              <[@{bond2}:]@{cl}\pcharge{-90}{Cl}
              >
              -
              ))
              -
              (<:[:-70])
              -
              (-[:30])
              (-[:-30])
              -@{carbon}
              =[@{dbond}:]
              )
            }
            \arrow{->}
            \chemfig[atomsep=1.5em]{
              [:-72]*5(
              =
              -
              (<[:-120]Cl)
              (<:[:-70])
              -
              (-[:30])
              (-[:-30])
              -
              (
              <[:90]@{oxygen}\pcharge{90}{O}
              (-[@{bond3}:15]@{hydrogen}H)
              (-[:165]H)
              )
              -
              )
            }
            \arrow{->[
              \chemfig{
                @{water2}\ce{H2O}
              }
            ]}
            \chemfig{
              [:-72]*5(
              =
              -
              (<[:-120]Cl)
              (<:[:-70])
              -
              (-[:30])
              (-[:-30])
              -
              (
              <[:90]OH
              )
              -
              )
            }
            \raisebox{-0.5cm}{
              \+
              \ce{H3O+}
              \+
              \ce{Cl-}
            }
          }
          \chemmovearrow{
            \reactionfigarrow{water1}{30}{carbon}{120}{Red}
            \reactionfigarrow{dbond}{-45}{bond1}{15}{Red}
            \reactionfigarrow{bond2}{-160}{cl}{-120}{Red}
            \reactionfigarrow{water2}{120}{hydrogen}{-60}{Red}
            \reactionfigarrow{bond3}{60}{oxygen}{100}{Red}
          }
        }{Secondary electrophilic carbon attack from front}{fig:ph29a}{0.75}

        \item \textbf{Secondary electrophilic carbon attack from back}

        \reactionfig{
          \scheme{
            \chemfig{
              @{water1}\ce{H2O}
            }
            \qquad \qquad
            \chemfig[atomsep=1.5em]{
              [:-72]*5(
              -[@{bond1}:]
              (*3(
              <[@{bond2}:]@{cl}\pcharge{-90}{Cl}
              >
              -
              ))
              -
              (<:[:-70])
              -
              (-[:30])
              (-[:-30])
              -@{carbon}
              =[@{dbond}:]
              )
            }
            \arrow{->}
            \chemfig[atomsep=1.5em]{
              [:-72]*5(
              =
              -
              (<[:-120]Cl)
              (<:[:-70])
              -
              (-[:30])
              (-[:-30])
              -
              (
              <:[:90]@{oxygen}\pcharge{90}{O}
              (-[@{bond3}:15]@{hydrogen}H)
              (-[:165]H)
              )
              -
              )
            }
            \arrow{->[
              \chemfig{
                @{water2}\ce{H2O}
              }
            ]}
            \chemfig{
              [:-72]*5(
              =
              -
              (<[:-120]Cl)
              (<:[:-70])
              -
              (-[:30])
              (-[:-30])
              -
              (
              <:[:90]OH
              )
              -
              )
            }
            \raisebox{-0.5cm}{
              \+
              \ce{H3O+}
              \+
              \ce{Cl-}
            }
          }
          \chemmovearrow{
            \reactionfigarrow{water1}{30}{carbon}{120}{Red}
            \reactionfigarrow{dbond}{-45}{bond1}{15}{Red}
            \reactionfigarrow{bond2}{-160}{cl}{-120}{Red}
            \reactionfigarrow{water2}{120}{hydrogen}{-60}{Red}
            \reactionfigarrow{bond3}{60}{oxygen}{100}{Red}
          }
        }{Secondary electrophilic carbon attack from back}{fig:ph30a}{0.75}

      \end{enumerate}

    \end{core}

    \begin{core}{Radical Bromination of Terminal Alkene}

      The radical bromination of a terminal alkene proceeds in a manner analogous to the reaction shown in \autoref{core:chlorination_example}, wherein resonance creates the possibility of multiple distinct modes of attack and, consequently, three distinct products.

      \ce{Br2} is first generated in situ through \textbf{N-Bromosuccinimide (NBS)}, which enables the selective production of \textbf{monobrominated products}. It should be noted that direct addition of pure \ce{Br2} to alkenes instead results in dihalogenation, as outlined in \autoref{alk-core:halogenation}.

      \reactionfig{
        \scheme{
          \chemname{
            \raisebox{-2cm}{
              \chemfig[atom sep=2em]{
                *5(-
                (=[:-70]O)
                -N
                (-[:0]Br)
                -
                (=[:70]O)
                --)
              }
            }
          }{N-Bromo-succinmide}
          \raisebox{-2cm}{
            \+
            \chemfig{
              \ce{HBr}
            }
          }
          \arrow{->}
          \chemfig[atom sep=2em]{
            *5(-
            (=[:-70]O)
            -N
            (-[:0]H)
            -
            (=[:70]O)
            --)
          }
          \+
          \chemfig{
            \ce{Br2}
          }
        }
      }{N-Bromo-succinmide allows to the slow release of \ce{Br2} to avoid alkenes adding \ce{Br2} directly}{fig:ph31a}{0.85}
      \vspace{-0.2cm}

      Following the spontaneous dissociation of \ce{Br2} into radical bromine atoms, the first propagation step occurs, generating a radical on the original alkene compound.

      \reactionfig{
        \scheme{
          \chemfig{
            [:0]@{br}\charge{0:4pt=\LARGE$\cdot$}{Br}
          }
          \quad
          \+
          \chemfig[atom sep=2em]{
            *5(-
            -@{carbon}
            (<[@{bond}:30]H)
            (
            <:[:-30]
            =[:30]
            )
            -
            (<[:50])
            (<:[:110]H)
            --)
          }
          \arrow{->}
          \chemfig[atom sep=2em]{
            *5(-
            -\charge{60:4pt=\LARGE$\cdot$}{}
            (
            -[:-30]
            =[:30]
            )
            -
            (<[:70])
            --)
          }
          \+
          \chemfig{
            [:0]HBr
          }
        }
        \chemmovearrow{
          \reactionfigarrow{bond}{90}{br}{60}{Green}
          \reactionfigarrow{bond}{-60}{carbon}{-70}{Green}
        }
      }{First propagation}{fig:ph32a}{0.9}

      This alkene radical and the remaining bromine radical can then undergo a second propagation step in three distinct manners.

      \begin{enumerate}

        \item \textbf{Propagation 2, Tertiary Radical Attack}

        \reactionfig{
          \scheme{
            \chemfig[atom sep=2em]{
              *5(-
              -@{radical}\charge{60:4pt=\LARGE$\cdot$}{}
              (
              -[:-30]
              =[:30]
              )
              -
              (<[:70])
              --)
            }
            \+
            \chemfig{
              [:0]Br
              -[@{bond}:0]@{br}Br
            }
            \arrow{->}
            \chemfig[atom sep=2em]{
              *5(-
              -
              (<:[:30]Br)
              (
              <[:-30]
              =[:30]
              )
              -
              (<[:70])
              --)
            }
            \+
            \chemfig[atom sep=2em]{
              *5(-
              -
              (<[:30]Br)
              (
              <:[:-30]
              =[:30]
              )
              -
              (<[:70])
              --)
            }
            \+
            \chemfig{
              \charge{0:4pt=\LARGE$\cdot$}{Br}
            }
          }
          \chemmovearrow{
            \reactionfigarrow{bond}{100}{br}{80}{Green}
            \reactionfigarrow{bond}{-100}{radical}{-75}{Green}
          }
        }{Propagation 2, tertiary radical attack}{fig:ph33a}{0.9}

        \item \textbf{Propagation 2, Terminal Radical Attack}

        \reactionfig{
          \scheme{
            \chemfig[atom sep=2em]{
              *5(-
              -@{radical}\charge{60:4pt=\LARGE$\cdot$}{}
              (
              -[@{bond1}:-30]
              =[@{bond2}:30]@{carbon}
              )
              -
              (<[:70])
              --)
            }
            \+
            \chemfig{
              [:0]Br
              -[@{bond3}:0]@{br}Br
            }
            \arrow{->}
            \chemfig[atom sep=2em]{
              *5(-
              -
              (
              =[:-30]
              -[:30]
              -[:-30]Br
              )
              -
              (<[:70])
              --)
            }
            \+
            \chemfig{
              \charge{0:4pt=\LARGE$\cdot$}{Br}
            }
          }
          \chemmovearrow{
            \reactionfigarrow{bond3}{100}{br}{80}{Green}
            \reactionfigarrow{bond3}{110}{carbon}{30}{Green}
            \reactionfigarrow{bond2}{120}{carbon}{90}{Green}
            \reactionfigarrow{bond2}{-60}{bond1}{-120}{Green}
            \reactionfigarrow{radical}{60}{bond1}{50}{Green}
          }
        }{Propagation 2, terminal radical attack}{fig:ph34a}{0.9}

        \item \textbf{Propagation 2, Rotated Terminal Radical Attack}

        \reactionfig{
          \scheme{
            \chemfig[atom sep=2em]{
              *5(-
              -\charge{60:4pt=\LARGE$\cdot$}{}
              (
              -[:-30,,,,Red,line width=2pt]
              =[:30]
              )
              -
              (<[:70])
              --)
            }
            \arrow{->[\textcolor{Red}{rotation}][]}
            \chemfig[atom sep=2em]{
              *5(-
              -@{radical}\charge{60:4pt=\LARGE$\cdot$}{}
              (
              -[@{bond1}:-30]
              =[@{bond2}:-90]@{carbon}
              )
              -
              (<[:70])
              --)
            }
            \+
            \chemfig{
              [:0]Br
              -[@{bond3}:0]@{br}Br
            }
            \arrow{->}
            \chemfig[atom sep=2em]{
              *5(-
              -
              (
              =[:-30]
              -[:-90]
              -[:-30]Br
              )
              -
              (<[:70])
              --)
            }
            \+
            \chemfig{
              \charge{0:4pt=\LARGE$\cdot$}{Br}
            }
          }
          \chemmovearrow{
            \reactionfigarrow{bond3}{100}{br}{80}{Green}
            \reactionfigarrow{bond3}{-135}{carbon}{-45}{Green}
            \reactionfigarrow{bond2}{15}{carbon}{-15}{Green}
            \reactionfigarrow{bond2}{-150}{bond1}{-150}{Green}
            \reactionfigarrow{radical}{60}{bond1}{50}{Green}
          }
        }{Propagation 2, rotated terminal radical attack}{fig:ph35a}{0.9}

      \end{enumerate}

      Taken together, the complete reaction proceeds as shown below.

      \reactionfig{
        \scheme{
          \chemfig[atom sep=2em]{
            *5(-
            -
            (
            <:[:0]
            =[:-60]
            )
            -
            (<[:70])
            --)
          }
          \raisebox{0.5cm}{
            \+
            \chemfig{
              NBS
            }
          }
          \qquad \qquad
          \arrow{->[hv][catalytic Br]}
          \qquad
          \chemfig{
            \text{\textit{\textbf{mono}brominated products}}
          }
        }
      }{Radical bromination of terminal alkene full reaction}{fig:ph36a}{0.9}

    \end{core}

